\section[Packing failures and Lipschitz images]{Packing failures and the Lipschitz-image conclusion}
\label{sec:packing-endpoint}

We now combine the two exceptional-family estimates with excess
propagation.  A flat descendant starts the propagation; a later failure
of bilateral flatness forces an oscillatory ancestor.  The finite counting
argument below converts this implication into a Carleson estimate without
requiring selected flat descendants to be disjoint.

Throughout this section, $\mu$ satisfies the hypotheses of
Theorem~\ref{thm:main}, $E=\supp\mu$ has positive diameter, and
$\mathcal D$ is the lattice of Proposition~\ref{prop:dyadic-lattice}.
Write $\mathcal D(R)=\{Q\in\mathcal D:Q\subset R\}$ and let $k(Q)$
be the generation of $Q$, so that $\ell(Q)=2^{-k(Q)}$.  Fix
\[
 H\geq 4(C_0+1),\qquad
 \beta_H(Q)=b_\mu(z_Q,H\ell(Q)).
\]
All constants below are uniform in the support center and scale.  Their
permitted dependencies are $d,n,C_{\rm AD},C_{\rm R}$ and the explicitly
chosen apertures and error thresholds.

\subsection{A finite seed-counting lemma}

The next lemma concerns only a nested measurable partition.  Its role is
to count parents by the mass of one selected descendant of each parent.

\begin{lemma}[Finite seed charging]\label{lem:seed-charging}
Let $R\in\mathcal D$ and let $\mathcal F$ be a finite family of
subcubes of $R$.  For each $Q\in\mathcal F$, choose a cube $S(Q)\subset Q$
such that, for fixed $N\in\mathbb N$ and $M\geq1$,
\[
 k(Q)\leq k(S(Q))\leq k(Q)+N,
 \qquad \mu(Q)\leq M\mu(S(Q)).
\]
Let $\mathcal O\subset\mathcal D(R)$ satisfy
$\sum_{T\in\mathcal O}\mu(T)\leq K\mu(R)$.
Suppose that whenever $Q,Q'\in\mathcal F$ satisfy
\[
 k(S(Q))<k(Q'),\qquad Q'\subset S(Q),
\]
there is a cube $T\in\mathcal O$ with
\[
 k(S(Q))\leq k(T)<k(Q'),\qquad Q'\subset T.
\]
Then
\[
 \sum_{Q\in\mathcal F}\mu(Q)
 \leq M(N+1)(1+K)\mu(R).
\]
\end{lemma}

\begin{proof}
For every eligible ordered pair $(Q,Q')$, choose one such $T$, and let
$\mathcal O_0$ be the resulting finite set.  Make these choices before
fixing a point.  Outside the common null set of lattice boundaries, fix
$x\in R$ and list the parents whose seeds contain $x$ in increasing order
of generation.  Their generations are distinct, because their parents
also contain $x$.

Starting with the first parent $Q_1$, group together every listed parent
of generation at most $k(S(Q_1))$.  This group has at most $N+1$ members.
If another parent remains, let $Q_2$ be the first one.  Nesting gives
$Q_2\subset S(Q_1)$, so the chosen charge for $(Q_1,Q_2)$ contains $x$
and has generation in
$[k(S(Q_1)),k(Q_2))$.  Repeat this procedure.  The intervals supplying
successive charges are disjoint: the next interval starts at
$k(S(Q_2))\geq k(Q_2)$.  Thus the number of groups is at most
$1+\sum_{T\in\mathcal O_0}\mathbf1_T(x)$, and
\[
 \sum_{Q\in\mathcal F}\mathbf1_{S(Q)}(x)
 \leq (N+1)\left(\mathbf1_R(x)
               +\sum_{T\in\mathcal O_0}\mathbf1_T(x)\right).
\]
Integrating this finite inequality and using the two mass bounds proves
the result.  In particular, overlapping seeds and arbitrarily long gaps
between selected generations cause no difficulty.
\end{proof}

\subsection{Propagation along a cube chain}

\begin{proposition}[Bilateral flatness has Carleson failures]
\label{prop:beta-packing}
For each $H\geq4(C_0+1)$ and $\varepsilon>0$, there is a constant
$C_{H,\varepsilon}<\infty$, depending only on
$d,n,C_{\rm AD},C_{\rm R},H,\varepsilon$, such that
\[
 \sum_{\substack{Q\subset R\\\beta_H(Q)\geq\varepsilon}}\mu(Q)
 \leq C_{H,\varepsilon}\mu(R)
 \qquad(R\in\mathcal D).
\]
\end{proposition}

\begin{proof}
We first fix all numerical parameters.  Apply
Theorem~\ref{thm:excess-propagation} with $\gamma=1/4$ and output
tolerance $\varepsilon$.  Write its resulting constants as
\[
 \delta=2^{-j_*},\qquad K_*=2j_*+3,\qquad A_*=8\cdot2^{K_*}.
\]
Choose
\[
 A_{\rm f}>H2^{K_*}+2C_0+1,\qquad
 0<\alpha<\min\left\{\frac{\varepsilon H}{4},
                  \frac{\delta H}{2\sqrt{G+1}}\right\},
 \qquad J_{\rm o}>A_*H+2C_0+1,
 \qquad v=\frac{\delta^3}{2}.
\]
Here $G\leq9^dC_{\rm AD}$ is supplied by
Lemma~\ref{lem:global-growth}.  Proposition~\ref{prop:oscillation-packing}
shows that
\[
 \mathcal O=\{T\in\mathcal D:
                  W_{J_{\rm o}}(T)>v\ell(T)^n\}
\]
is Carleson; denote its constant by $C_{\rm o}$.
Proposition~\ref{prop:flat-seeds}, with aperture $A_{\rm f}$ and
tolerance $\alpha/(2A_{\rm f})$, supplies a Carleson family
$\mathcal C$, a fixed depth $N$, and, for every $Q\notin\mathcal C$,
a descendant $S(Q)\subset Q$ at depth at most $N$ with
\[
 b_\mu(z_{S(Q)},A_{\rm f}\ell(S(Q)))
       <\frac{\alpha}{2A_{\rm f}}.
\]
A near-minimizing plane $L_S$ consequently satisfies both directed
bounds
\begin{equation}\label{eq:seed-plane-bounds}
 \sup_{E\cap B(z_S,A_{\rm f}\ell(S))}\dist(\cdot,L_S)
 \leq\alpha\ell(S),\qquad
 \sup_{L_S\cap B(z_S,A_{\rm f}\ell(S))}\dist(\cdot,E)
 \leq\alpha\ell(S).
\end{equation}
The lattice mass bounds give a uniform $M\geq1$ such that
\begin{equation}\label{eq:seed-mass-comparison}
 \mu(Q)\leq M\mu(S(Q)),\qquad
 M\leq C(d,n,C_{\rm AD})2^{nN}.
\end{equation}

Before iterating propagation, we ensure its admissible horizon.
The required condition is
$2^{K_*}H\ell(S)\leq D_E$. If $D_E<\infty$, restrict for now to
parents satisfying $\ell(Q)\leq D_E/B_*$, where
$B_*=H2^{K_*}$. Every selected seed then satisfies the condition,
as do all later radii in its chain. If $D_E=\infty$, there is no cutoff.
This restriction costs only finitely many generations:
the largest lattice scale lies in $[D_E,2D_E)$, so there are at most
$2+\lceil\log_2 B_*\rceil$ omitted generations. At each generation,
subcubes of a given $R$ are disjoint and have total mass at most $\mu(R)$.
We will restore all omitted cubes at this fixed Carleson cost.

We prove the implication needed by Lemma~\ref{lem:seed-charging}.
Let $S$ be one of these flat seeds and let $Q'\subset S$ be a proper
descendant, with
$m=k(Q')-k(S)>0$.  Suppose that every proper ancestor of $Q'$ between
$S$ and $Q'$ lies outside $\mathcal O$.  We claim that
$\beta_H(Q')<\varepsilon$.

Keep the center $a=z_{Q'}$ fixed throughout the argument.  Let $T_j$
be the ancestor of $Q'$ of generation $k(S)+j$, and set
\[
 r_j=H2^{-j}\ell(S)=H\ell(T_j),\qquad 0\leq j\leq m.
\]
The center belongs to the closure of every $T_j$, so
$|a-z_{T_j}|\leq2C_0\ell(T_j)$.  In particular,
$B(a,2^{K_*}r_0)\subset B(z_S,A_{\rm f}\ell(S))$.
The first bound in \eqref{eq:seed-plane-bounds} and global upper growth
give, for $0\leq i\leq K_*$,
\begin{equation}\label{eq:seed-initial-excess}
 e_\mu(a,2^ir_0)
 \leq\frac{\sqrt G\,\alpha\ell(S)}{2^ir_0}
 \leq\delta\,2^{i/4}.
\end{equation}

For $j<m$, a compactly supported mean-zero test $\varphi$ with
$\Lip\varphi\leq r_j^{-1}$ and
$\supp\varphi\subset B(a,A_*r_j)$ is admissible in the definition of
$W_{J_{\rm o}}(T_j)$: its support lies in the required enlarged ball,
and $r_j\geq\ell(T_j)$.  Therefore
\begin{equation}\label{eq:chain-oscillation}
 |\mathcal R_\mu(\varphi)|
 \leq v\ell(T_j)^n\leq\delta^3r_j^n.
\end{equation}
We may now iterate Theorem~\ref{thm:excess-propagation} at the radii
$r_0,r_1,\ldots,r_{m-1}$.  To check the complete horizon at the $j$th
step, observe that
\[
 2^ir_j=
 \begin{cases}
 r_{j-i},&i\leq j,\\
 2^{i-j}r_0,&i>j.
 \end{cases}
\]
In the first case the preceding induction gives excess at most $\delta$;
in the second, \eqref{eq:seed-initial-excess} gives at most
$\delta2^{(i-j)/4}\leq\delta2^{i/4}$.  Thus all horizon inequalities
hold at every step.  Together with \eqref{eq:chain-oscillation}, they
prove
\[
 e_\mu(a,r_j)\leq\delta\quad(0\leq j\leq m),\qquad
 b_\mu(a,r_j)<\varepsilon\quad(1\leq j\leq m).
\]
Since $r_m=H\ell(Q')$, the claim follows.  Notice that no test at the
terminal cube $Q'$ was used.

Fix $R\in\mathcal D$.  Take any finite family $\mathcal F$ of
remaining subcubes of $R$ with $\beta_H(Q)\geq\varepsilon$, and remove
the members of $\mathcal C$.
For each remaining parent use the seed selected above.  If
$Q'\subset S(Q)$ and $k(Q')>k(S(Q))$, then the just-proved implication,
applied contrapositively, supplies $T\in\mathcal O$ with
\[
 S(Q)\supset T\supset Q',\qquad
 k(S(Q))\leq k(T)<k(Q').
\]
Lemma~\ref{lem:seed-charging} and
\eqref{eq:seed-mass-comparison} give
\[
 \sum_{Q\in\mathcal F\setminus\mathcal C}\mu(Q)
 \leq M(N+1)(1+C_{\rm o})\mu(R).
\]
Add the Carleson bound for $\mathcal C$ and the finite coarse-generation
cost.  Taking the supremum over finite families proves the asserted
countable sum.  Every parameter was fixed before $R$ or any seed was
chosen.
\end{proof}

\subsection{From dyadic packing to the bilateral weak geometric lemma}

We spell out the passage to the intrinsic ball condition, so the final
geometric theorem does not depend on a choice of lattice.  First,
$b_\mu$ is Borel on $E\times(0,\infty)$.  Indeed, the two unnormalized
directed deviations on an open ball are jointly lower semicontinuous
in its center, radius, and affine plane.  Points strictly inside a ball
persist under small perturbations, and points of a limiting plane can
be approximated by points of the varying planes.  Along a convergent
sequence of centers and positive radii, near-minimizing planes with
bounded error have bounded offsets: the center belongs to $E$ and its
distance to the plane is one of the first directed deviations.
Compactness of the Grassmannian then supplies a convergent plane
subsequence.  The two lower-semicontinuity inequalities show that their
infimum $b_\mu$ is lower semicontinuous as well.

Fix $\varepsilon>0$, $a\in E$, and $0<R<D_E$.  Choose a lattice scale
$s$ with $R\leq s<2R$, and put $s_j=2^{-j}s$.  For almost every $x$,
let $Q_j(x)$ be its cube of side scale $s_j$.  Whenever
$s_j/2<t\leq s_j$,
\[
 B(x,t)\subset B(z_{Q_j(x)},Hs_j),\qquad
 b_\mu(x,t)\leq\frac{Hs_j}{t}\beta_H(Q_j(x))
                  \leq2H\beta_H(Q_j(x)).
\]
Both inequalities follow directly from the two directed deviations;
the same plane is used before taking the infimum.  Integrating over
the dyadic scale intervals and then over $E\cap B(a,R)$ yields
\begin{align*}
 &\int_{E\cap B(a,R)}\int_0^R
       \mathbf1_{\{b_\mu(x,t)>\varepsilon\}}\,\frac{dt}{t}\,d\mu(x)\\
 &\qquad\leq(\log2)
   \sum_{\substack{Q\text{ at scale }s_j,\ j\geq0\\
          Q\cap B(a,R)\ne\varnothing\\
          \beta_H(Q)>\varepsilon/(2H)}}\mu(Q).
\end{align*}
Group these cubes by their scale-$s$ ancestors $P$.  Such a $P$ meets
$B(a,R)$ and is contained in $B(a,(1+4C_0)R)$.  Applying
Proposition~\ref{prop:beta-packing} to each $P$, then summing their
disjoint masses and using global growth, proves
\begin{equation}\label{eq:continuous-bwgl}
 \int_{E\cap B(a,R)}\int_0^R
       \mathbf1_{\{b_\mu(x,t)>\varepsilon\}}\,\frac{dt}{t}\,d\mu(x)
 \leq C_\varepsilon R^n.
\end{equation}
This is the \emph{bilateral weak geometric lemma}: for every positive
flatness threshold, its failures occupy a Carleson set of positions
and scales.

\subsection{The geometric criterion and the ball-map conclusion}

For clarity, we recall the measure comparison used to apply the
geometric criterion.  AD regularity with $D_E>0$ gives
\begin{equation}\label{eq:measure-hausdorff-comparison}
 c\,\H^n\!\restriction E\ \leq\ \mu\ \leq\
 C\,\H^n\!\restriction E,
\end{equation}
with positive finite constants depending only on the AD data and the
Hausdorff-measure normalization.  Here is a direct justification.
Cover a subset of $E$ by sets of arbitrarily small diameter, recenter
each nonempty set at a point of $E$, and use the upper ball bound.
Infimizing the cover sums proves the upper inequality.  For the reverse
inequality, let $F\subset E$ be bounded and Borel, and let $U$ be an open
set containing $F$.  Cover $F$ by support-centered balls contained in
$U$, of radii less than both $D_E/10$ and an arbitrarily small prescribed
number.  The $5r$ covering lemma selects disjoint balls $B_i$ whose
fivefold dilates cover $F$.  The lower AD bound gives
\[
 \H^n_\eta(F)\leq C_n\sum_i r_i^n
      \leq C_nC_{\rm AD}\sum_i\mu(B_i)
      \leq C_nC_{\rm AD}\mu(U),
\]
Here $\H^n_\eta$ denotes Hausdorff content using covers of diameter
less than $\eta$, with the same normalization as $\H^n$; the chosen
radii make all fivefold diameters less than $\eta$.
Let $\eta\downarrow0$, use outer regularity, and then exhaust unbounded
sets by bounded sets.  This proves \eqref{eq:measure-hausdorff-comparison}.
In particular $E$ is an $n$-AD-regular set for $\H^n$.  Write $b_E$ for
the bilateral coefficient of Definition~\ref{def:flatness} with support
set $E$, so that $b_E=b_\mu$.  Replacing
$d\mu$ by $d\H^n$ in \eqref{eq:continuous-bwgl} changes only its constant.

\begin{theorem}[David--Semmes bilateral criterion]\label{thm:bwgl}
Let $E\subset\R^d$ be closed, of positive diameter, and $n$-AD regular
with respect to $\H^n$.  Suppose that for every $\varepsilon>0$ there is
$C_\varepsilon<\infty$ such that
\[
 \int_{E\cap B(a,R)}\int_0^R
   \mathbf1_{\{b_E(x,t)>\varepsilon\}}\,\frac{dt}{t}\,d\H^n(x)
 \leq C_\varepsilon R^n
 \quad(a\in E,\ 0<R<\diam E).
\]
Then $E$ has big pieces of uniformly Lipschitz images: there are
$\theta_0>0$ and $L<\infty$ such that for every $a\in E$ and
$0<r<\diam E$ there is an $L$-Lipschitz map
$f:B^n(0,r)\to\R^d$ for which
\[
 \H^n\bigl(E\cap B(a,r)\cap f(B^n(0,r))\bigr)\geq\theta_0r^n.
\]
The constants depend only on the dimensions, AD constants, and the
constants in the bilateral weak geometric lemma.
\end{theorem}

This is the Euclidean David--Semmes criterion
\cite[Chapter~II.2]{DavidSemmes1993}; see also
\cite[Theorem~2.5]{Tolsa2015Uniform} for the all-scale formulation and
\cite[Theorem~1.0.4]{BateHydeSchul2023} for finite diameter.
Appendix~\ref{app:geometric-conventions} verifies the finite-diameter,
open-ball and positive-mass conventions needed for this formulation.

\begin{proof}[Proof of Theorem~\ref{thm:main}]
If $E$ is empty or a singleton, there is no positive admissible radius,
and the asserted ball-map conclusion is vacuous.  Otherwise
Proposition~\ref{prop:beta-packing} gives
\eqref{eq:continuous-bwgl} for every threshold.  The comparison
\eqref{eq:measure-hausdorff-comparison} permits application of
Theorem~\ref{thm:bwgl}; its positive image-mass bound also holds with
$\mu$, after changing the positive constant.

For any $a\in E$ and $0<r\leq D_E$, apply this conclusion at $r/2$ and
write its map as $f:B^n(0,r/2)\to\R^d$.  The map
\[
 g:B^n(0,r)\longrightarrow\R^d,\qquad g(u)=f(u/2),
\]
has the same image and Lipschitz constant at most $L/2$.  Since
$B(a,r/2)\subset B(a,r)$, it captures at least
$c\theta_0 2^{-n}r^n$ of the original measure in $B(a,r)$.
This also covers $r=D_E$ when the diameter is finite.  All constants
were chosen before $a$ and $r$, and have only the dependencies asserted
in the theorem.  Thus the uniform ball-map condition in
Definition~\ref{def:ur} holds.
\end{proof}
